% ========================================================
%  习近平新时代中国特色社会主义思想概论  全真模拟仿真试卷（含参考答案）
%  覆盖范围：导论、第一章～第十二章（依据课程复习纲要）
%  试卷结构：30 道单项选择（30 分）+ 10 道多项选择（20 分）+ 5 道简答题（50 分）
% ========================================================
\documentclass[12pt, a4paper]{article}

% ---- 中文支持 ----
\usepackage[UTF8]{ctex}

% ---- 数学 ----
\usepackage{amsmath, amssymb, amsthm}
\usepackage{mathtools}

% ---- 页面布局 ----
\usepackage[top=2.5cm, bottom=2.5cm, left=2.8cm, right=2.8cm]{geometry}

% ---- 表格 ----
\usepackage{booktabs, array, multirow, makecell}
\usepackage{tabularx}

% ---- 页眉页脚 ----
\usepackage{fancyhdr}
\usepackage{lastpage}
\setlength{\headheight}{14pt}
\pagestyle{fancy}
\fancyhf{}
\renewcommand{\headrulewidth}{0.6pt}
\lhead{\small 《习概》全真模拟仿真试卷·MiniMax卷}
\rhead{\small 共 \pageref{LastPage} 页}
\cfoot{\small 第 \thepage\ 页}

% ---- 计数器与样式 ----
\usepackage{enumitem}
\usepackage{xcolor}
\usepackage{tcolorbox}
\tcbuselibrary{skins, breakable}

% ---- 填空线 ----
\newcommand{\blank}[1]{\underline{\hspace{#1}}}

% ---- 答案盒子 ----
\newtcolorbox{answerbox}[1][]{
  breakable,
  colback=gray!6,
  colframe=gray!50,
  fonttitle=\bfseries\small,
  title=参考答案,
  #1
}

% ---- 题号格式 ----
\newcounter{bigq}
\newcommand{\BigQ}[1]{%
  \stepcounter{bigq}%
  \vspace{6pt}%
  \noindent{\large\bfseries 第\chinese{bigq}大题\quad #1}%
  \vspace{4pt}\par
}

\newcounter{subq}[bigq]
\newcommand{\SubQ}{%
  \stepcounter{subq}%
  \par\noindent\textbf{（\arabic{subq}）}\quad
}

% 分隔线
\newcommand{\examrule}{\noindent\rule{\linewidth}{0.6pt}}

% 选项排版辅助宏（A、B 在一行；C、D 在下一行）
\newcommand{\opts}[4]{%
A. #1 \quad B. #2\\[2pt]
C. #3 \quad D. #4}

% ---- 文档开始 ----
\begin{document}

% ============================================================
%  封面
% ============================================================
\begin{center}
  {\LARGE\bfseries 《习近平新时代中国特色社会主义思想概论》}\\[8pt]
  {\large\bfseries 全真模拟仿真试卷}\\[16pt]
  \begin{tabular}{lll}
    考试时间：120 分钟 & \quad & 满分：100 分 \\[4pt]
    \multicolumn{3}{l}{注意事项：请将答案写在答题纸指定区域，写在本卷上无效。多选题少选、错选均不得分。}
  \end{tabular}
  \vspace{4pt}

  \examrule

  \vspace{6pt}
  \begin{tabular}{|c|c|c|c|c|}
    \hline
    \textbf{题号} & 一 & 二 & 三 & \textbf{总分} \\
    \hline
    \textbf{满分} & 30 & 20 & 50 & 100 \\
    \hline
    \textbf{得分} & \quad & \quad & \quad & \quad \\
    \hline
  \end{tabular}
\end{center}

\vspace{10pt}
\examrule
\vspace{6pt}

% ============================================================
%  第一大题：单项选择题
% ============================================================
\BigQ{单项选择题（每题 1 分，共 30 分）}

\SubQ 习近平新时代中国特色社会主义思想创立的时代背景中，世界正经历\blank{2em}，国际环境日趋复杂。
\opts{百年未有之大变局加速演进}{冷战格局全面形成}{多极化格局彻底瓦解}{全球化进程全面逆转}
\SubQ 在"两个结合"中，马克思主义基本原理被称为"\blank{2em}"，中华优秀传统文化被称为"\blank{2em}"。
\opts{根脉；魂脉}{魂脉；根脉}{基础；核心}{核心；基础}
\SubQ 下列关于习近平新时代中国特色社会主义思想历史地位的表述中，正确的是\blank{2em}。
\opts{马克思主义中国化时代化的第三次飞跃}{当代中国马克思主义、二十一世纪马克思主义}{彻底否定了马克思主义的普遍真理}{仅适用于中国特色社会主义初级阶段}
\SubQ 习近平新时代中国特色社会主义思想科学体系中，起"四梁八柱"支撑作用的核心组成部分是\blank{2em}。
\opts{十个明确}{十四个坚持}{十三个方面成就}{六个必须坚持}
\SubQ 中国特色社会主义进入新时代，我国社会主要矛盾已经转化为\blank{2em}。
\opts{人民日益增长的物质文化需要同落后的社会生产之间的矛盾}{人民日益增长的美好生活需要和不平衡不充分的发展之间的矛盾}{人民日益增长的精神文化需要同不平衡不充分的发展之间的矛盾}{人民日益增长的美好生活需要和经济发展方式之间的矛盾}
\SubQ "一个变、两个没有变"中，没有变的是\blank{2em}。
\opts{社会主要矛盾}{我国仍处于并将长期处于社会主义初级阶段的基本国情}{国际格局}{党的中心任务}
\SubQ "四个全面"战略布局中，\blank{2em}是战略目标。
\opts{全面深化改革}{全面依法治国}{全面建设社会主义现代化国家}{全面从严治党}
\SubQ 中国式现代化的本质要求中，\blank{2em}是政治保证。
\opts{坚持中国特色社会主义}{坚持中国共产党领导}{实现高质量发展}{推动构建人类命运共同体}
\SubQ 中国式现代化的中国特色中，下列说法错误的是\blank{2em}。
\opts{人口规模巨大的现代化}{全体人民共同富裕的现代化}{物质文明和精神文明相协调的现代化}{以资本为中心的现代化}
\SubQ 推进中国式现代化需要牢牢把握的重大原则中，\blank{2em}处于首位。
\opts{坚持中国特色社会主义道路}{坚持以人民为中心的发展思想}{坚持和加强党的全面领导}{坚持发扬斗争精神}
\SubQ 中国最大的国情是\blank{2em}。
\opts{处于并将长期处于社会主义初级阶段}{中国共产党的领导}{是世界上最大的发展中国家}{实行社会主义市场经济}
\SubQ 党的自身优势中，\blank{2em}是马克思主义政党第一位的能力。
\opts{思想引领力}{群众组织力}{政治领导力}{社会号召力}
\SubQ 关于党的领导，下列表述正确的是\blank{2em}。
\opts{党的领导意味着事事都要由党具体包办}{党的领导是局部的、有条件的}{党的领导是全面的、系统的、整体的}{党的领导只对党内有效}
\SubQ 中国共产党的根本政治立场是\blank{2em}。
\opts{阶级立场}{国家立场}{人民立场}{民族立场}
\SubQ "江山就是人民，人民就是江山"中，"江山"指的是\blank{2em}。
\opts{祖国的山河}{党的执政地位}{人民的幸福生活}{民族复兴的大业}
\SubQ 共同富裕\blank{2em}整齐划一的同步实现。
\opts{是}{不是}{等同于}{必须}
\SubQ 党的\blank{2em}是国家的生命线、人民的幸福线。
\opts{基本理论}{基本路线}{基本方略}{基本纲领}
\SubQ 全面深化改革的总目标是\blank{2em}。
\opts{实现中华民族伟大复兴}{完善和发展中国特色社会主义制度、推进国家治理体系和治理能力现代化}{建成社会主义现代化强国}{实现共同富裕}
\SubQ 党的十八届三中全会通过了\blank{2em}，对全面深化改革进行了总体擘画。
\opts{《关于建国以来党的若干历史问题的决议》}{《中共中央关于全面深化改革若干重大问题的决定》}{《中共中央关于党的百年奋斗重大成就和历史经验的决议》}{《关于新形势下党内政治生活的若干准则》}
\SubQ "三个进一步解放"不包括\blank{2em}。
\opts{进一步解放思想}{进一步解放和发展社会生产力}{进一步解放和增强社会活力}{进一步解放和扩大对外开放}
\SubQ 新发展理念中，\blank{2em}是国家繁荣发展的必由之路。
\opts{创新}{协调}{绿色}{开放}
\SubQ 我国经济已由\blank{2em}转向\blank{2em}。
\opts{高速增长阶段；中高速增长阶段}{高速增长阶段；高质量发展阶段}{数量型增长阶段；质量型增长阶段}{要素驱动阶段；创新驱动阶段}
\SubQ 新发展格局的核心要义是\blank{2em}。
\opts{以国际大循环为主体}{以国内大循环为主体、国内国际双循环相互促进}{完全依赖国内市场}{出口主导、进口为辅}
\SubQ "三个第一"指的是\blank{2em}。
\opts{科技是第一生产力、人才是第一资源、创新是第一动力}{教育是第一基础、人才是第一资本、创新是第一引擎}{科技是第一竞争力、人才是第一财富、创新是第一动力}{教育是根本、科技是关键、人才是基础}
\SubQ 教育、科技、人才在全面建设社会主义现代化国家中的定位是\blank{2em}支撑。
\opts{基础性、关键性、先导性}{基础性、战略性、先导性}{基础性、战略性、支撑性}{关键性、支撑性、先导性}
\SubQ 人民民主是\blank{2em}的生命。
\opts{中华民族}{中国共产党}{社会主义}{中国特色社会主义制度}
\SubQ 在马克思主义看来，民主首先是、而且主要是一种\blank{2em}。
\opts{思想观念}{价值追求}{国家制度}{生活方式}
\SubQ 中国特色社会主义文化的三个主要组成部分中不包括\blank{2em}。
\opts{中华优秀传统文化}{革命文化}{社会主义先进文化}{西方资本主义文化}
\SubQ 加快建设\blank{2em}是化解风险矛盾最直接、最有效率的治理主体。
\opts{基层社会治理现代化}{社区治理现代化}{市域社会治理现代化}{乡村治理现代化}
\SubQ 生态文明建设应摆在\blank{2em}。
\opts{经济发展的优先位置}{全局工作的突出位置}{现代化建设的关键位置}{城市发展的核心位置}

\vspace{10pt}
\examrule

% ============================================================
%  第二大题：多项选择题
% ============================================================
\BigQ{多项选择题（每题 2 分，共 20 分；每题至少 2 项正确，多选、错选、少选均不得分）}

\SubQ 习近平新时代中国特色社会主义思想是"两个结合"的重大成果，"两个结合"具体指\blank{2em}。
\opts{马克思主义基本原理同中国具体实际相结合}{马克思主义基本原理同中华优秀传统文化相结合}{马克思主义基本原理同中国革命实践相结合}{马克思主义基本原理同西方现代文明相结合}
\SubQ 中国特色社会主义进入新时代，这个新时代是\blank{2em}的时代。
\opts{承前启后、继往开来、在新的历史条件下继续夺取中国特色社会主义伟大胜利}{决胜全面建成小康社会、进而全面建设社会主义现代化强国}{全国各族人民团结奋斗、不断创造美好生活、逐步实现全体人民共同富裕}{全体中华儿女勠力同心、奋力实现中华民族伟大复兴中国梦}
\SubQ 中国式现代化的本质要求包括\blank{2em}。
\opts{坚持中国共产党领导}{坚持中国特色社会主义}{实现高质量发展}{推动构建人类命运共同体}
\SubQ 中国共产党的自身优势主要包括\blank{2em}。
\opts{政治领导力}{思想引领力}{群众组织力}{社会号召力}
\SubQ 坚持以人民为中心的发展思想，要求做到\blank{2em}。
\opts{发展为了人民}{发展依靠人民}{发展成果由人民共享}{把人民对美好生活的向往作为奋斗目标}
\SubQ 新时代全面深化改革开放是一场深刻革命，主要体现在\blank{2em}。
\opts{是一场思想理论的深刻变革}{是改革组织方式的深刻变革}{是国家制度和治理体系的深刻变革}{是人民广泛参与的深刻变革}
\SubQ 新发展理念包括\blank{2em}。
A. 创新 \quad B. 协调 \quad C. 绿色\\[2pt]
D. 开放 \quad E. 共享
\SubQ 全面建设社会主义现代化国家要实施的三大战略是\blank{2em}。
\opts{科教兴国战略}{人才强国战略}{创新驱动发展战略}{区域协调发展战略}
\SubQ 全过程人民民主体现的"四个相统一"包括\blank{2em}。
\opts{过程民主和成果民主相统一}{程序民主和实质民主相统一}{直接民主和间接民主相统一}{人民民主和国家意志相统一}
\SubQ 关于文化、民生与生态文明建设，下列表述正确的有\blank{2em}。
\opts{文化繁荣兴盛是实现中华民族伟大复兴的必然要求}{让人民生活幸福是"国之大者"}{生态文明建设应摆在全局工作的突出位置}{意识形态关乎旗帜、关乎道路、关乎国家安全}

\vspace{10pt}
\examrule

% ============================================================
%  第三大题：简答题
% ============================================================
\BigQ{简答题（每题 10 分，共 50 分；要求观点正确、要点完整，并进行简要论述）}

\SubQ（10 分）\textbf{简述习近平新时代中国特色社会主义思想创立的时代背景，并论述"两个确立"的决定性意义。}

\vspace{2.5em}

\SubQ（10 分）\textbf{简述中国式现代化的中国特色，并论述推进中国式现代化需要牢牢把握的重大原则。}

\vspace{2.5em}

\SubQ（10 分）\textbf{为什么说中国共产党的领导是中国特色社会主义最本质的特征？新时代如何健全党的全面领导制度？}

\vspace{2.5em}

\SubQ（10 分）\textbf{为什么说人民立场是中国共产党的根本政治立场？如何理解共同富裕是中国特色社会主义的本质要求？}

\vspace{2.5em}

\SubQ（10 分）\textbf{为什么说新时代全面深化改革开放是一场深刻革命？结合新发展理念，谈谈如何以新发展理念引领高质量发展。}

\vspace{2.5em}

\vspace{10pt}
\examrule

\vspace{16pt}
\examrule
\vspace{4pt}
\begin{center}
  {\large\bfseries ——试题到此结束，请认真检查——}
\end{center}
\vspace{4pt}
\examrule

% ============================================================
%  参考答案
% ============================================================
\newpage
\setcounter{bigq}{0}

\begin{center}
  {\LARGE\bfseries 参考答案与评分标准}
\end{center}
\vspace{8pt}
\examrule

% -------- 第一大题参考答案 --------
\BigQ{单项选择题参考答案（每题 1 分，共 30 分）}

\begin{answerbox}
\begin{tabular}{cl@{\hspace{1.2em}}cl@{\hspace{1.2em}}cl@{\hspace{1.2em}}cl}
1.\ A & 2.\ B & 3.\ B & 4.\ A \\
5.\ B & 6.\ B & 7.\ C & 8.\ B \\
9.\ D & 10.\ C & 11.\ B & 12.\ C \\
13.\ C & 14.\ C & 15.\ B & 16.\ B \\
17.\ B & 18.\ B & 19.\ B & 20.\ D \\
21.\ D & 22.\ B & 23.\ B & 24.\ A \\
25.\ C & 26.\ C & 27.\ C & 28.\ D \\
29.\ C & 30.\ B & & \\
\end{tabular}
\end{answerbox}

% -------- 第二大题参考答案 --------
\BigQ{多项选择题参考答案（每题 2 分，共 20 分）}

\begin{answerbox}
\begin{tabular}{cl@{\hspace{1.2em}}cl@{\hspace{1.2em}}cl@{\hspace{1.2em}}cl}
1.\ AB & 2.\ ABCD & 3.\ ABCD & 4.\ ABCD \\
5.\ ABCD & 6.\ ABCD & 7.\ ABCDE & 8.\ ABC \\
9.\ ABCD & 10.\ ABCD & & \\
\end{tabular}
\end{answerbox}

% -------- 第三大题参考答案 --------
\BigQ{简答题参考答案（共 50 分）}

\begin{answerbox}
\noindent\textbf{第 1 题（10 分）：简述习近平新时代中国特色社会主义思想创立的时代背景，并论述"两个确立"的决定性意义。}

\vspace{4pt}
\textbf{【参考答案】}

\noindent（一）时代背景（5 分）：

（1）世界百年未有之大变局加速演进，国际力量对比深刻调整，国际环境日趋复杂，不稳定性不确定性明显增加。

（2）中华民族伟大复兴进入关键时期，改革发展稳定任务艰巨繁重。

（3）中国式现代化全面推进拓展，社会主义现代化建设进入新阶段。

（4）科学社会主义在 21 世纪的中国焕发新的蓬勃生机，中国特色社会主义展现出强大生命力。

（5）中国共产党在革命性锻造中更加坚强有力，党的执政能力和领导水平不断提高。

\noindent（二）"两个确立"的决定性意义（5 分）：

（1）"两个确立"指：确立习近平同志党中央的核心、全党的核心地位；确立习近平新时代中国特色社会主义思想的指导地位。

（2）"两个确立"是新时代取得一切成就的根本原因，是战胜一切艰难险阻、应对一切不确定性的最大确定性、最大底气、最大保证。

（3）"两个确立"反映了全党全军全国各族人民共同心愿，对新时代党和国家事业发展、对推进中华民族伟大复兴历史进程具有决定性意义。

\vspace{6pt}
\noindent\rule{\linewidth}{0.4pt}
\vspace{4pt}

\noindent\textbf{第 2 题（10 分）：简述中国式现代化的中国特色，并论述推进中国式现代化需要牢牢把握的重大原则。}

\vspace{4pt}
\textbf{【参考答案】}

\noindent（一）中国式现代化的中国特色（5 分）：

（1）人口规模巨大的现代化——14 亿多人口整体迈进现代化社会，规模超过现有发达国家人口的总和。

（2）全体人民共同富裕的现代化——坚持以人民为中心的发展思想，避免两极分化。

（3）物质文明和精神文明相协调的现代化——既物质富足又精神富足，促进物的全面丰富和人的全面发展。

（4）人与自然和谐共生的现代化——坚持可持续发展，统筹推进经济建设和生态环境保护。

（5）走和平发展道路的现代化——通过合作共赢、互利互惠拓展发展空间，坚决不走国强必霸的老路。

\noindent（二）推进中国式现代化需要牢牢把握的重大原则（5 分）：

（1）坚持和加强党的全面领导——确保党始终成为坚强领导核心。

（2）坚持中国特色社会主义道路——保持方向不变、道路不偏。

（3）坚持以人民为中心的发展思想——让现代化建设成果更多更公平惠及全体人民。

（4）坚持深化改革开放——为现代化建设注入动力和活力。

（5）坚持发扬斗争精神——敢于斗争、善于斗争，应对前进路上的风险挑战。

\vspace{6pt}
\noindent\rule{\linewidth}{0.4pt}
\vspace{4pt}

\noindent\textbf{第 3 题（10 分）：为什么说中国共产党的领导是中国特色社会主义最本质的特征？新时代如何健全党的全面领导制度？}

\vspace{4pt}
\textbf{【参考答案】}

\noindent（一）党的领导是中国特色社会主义最本质的特征（5 分）：

（1）中国最大的国情就是中国共产党的领导，这是中国特色社会主义最本质的特征，也是中国特色社会主义制度的最大优势。

（2）这一特征是由科学社会主义的理论逻辑、中国特色社会主义产生发展的历史逻辑、中国共产党领导中国革命建设改革的实践逻辑共同决定的。

（3）中国共产党的自身优势（政治领导力、思想引领力、群众组织力、社会号召力）是中国特色社会主义制度优势的主要来源，党的领导是党和国家事业不断发展的"定海神针"。

\noindent（二）健全党的全面领导制度（5 分）：

（1）必须完善党在各种组织中发挥领导作用的制度——在人大、政府、政协、监察机关、司法机关以及人民团体、企事业单位、基层群众性自治组织、社会组织中，党发挥领导核心作用。

（2）必须完善党协调各方的机制——统筹协调人大、政府、政协、纪委监委、法院、检察院、人民团体、企事业单位以及社会组织等各方力量，形成工作合力。

（3）必须完善党领导各项事业的具体制度——把党的领导贯彻到党和国家所有机构履行职责全过程、各方面，实现党的领导全覆盖、全贯穿。

\vspace{6pt}
\noindent\rule{\linewidth}{0.4pt}
\vspace{4pt}

\noindent\textbf{第 4 题（10 分）：为什么说人民立场是中国共产党的根本政治立场？如何理解共同富裕是中国特色社会主义的本质要求？}

\vspace{4pt}
\textbf{【参考答案】}

\noindent（一）人民立场是中国共产党的根本政治立场（5 分）：

（1）人民性是马克思主义的本质属性和鲜明品格，中国共产党来自人民、依靠人民，党的一切奋斗都是为了人民。

（2）坚持人民立场，就要始终牢记党的初心和使命，把为人民谋幸福作为党的根本使命。

（3）坚持人民立场，就要始终保持党同人民群众的血肉联系，热爱人民、尊重人民、敬畏人民。

（4）民心是最大的政治，决定事业兴衰成败；人民是党的工作的最高裁决者和最终评判者，"时代是出卷人，我们是答卷人，人民是阅卷人"。

\noindent（二）共同富裕是中国特色社会主义的本质要求（5 分）：

（1）共同富裕是中国特色社会主义的根本原则，是中国式现代化的重要特征，关系到党的执政基础的重大政治问题。

（2）共同富裕不是整齐划一的同步实现、同等富裕，必须遵循"鼓励勤劳创新致富、坚持基本经济制度、尽力而为量力而行、坚持循序渐进"的原则。

（3）要在高质量发展中促进共同富裕，正确处理效率和公平的关系，构建初次分配、再分配、三次分配协调配套的基础性制度安排，形成中间大、两头小的橄榄型分配结构。

（4）推动全体人民共同富裕与促进人的全面发展是高度统一的，要让发展成果更多更公平惠及全体人民。

\vspace{6pt}
\noindent\rule{\linewidth}{0.4pt}
\vspace{4pt}

\noindent\textbf{第 5 题（10 分）：为什么说新时代全面深化改革开放是一场深刻革命？结合新发展理念，谈谈如何以新发展理念引领高质量发展。}

\vspace{4pt}
\textbf{【参考答案】}

\noindent（一）新时代全面深化改革开放是一场深刻革命（5 分）：

（1）就其艰巨性、复杂性和系统性来说，全面深化改革开放是一场涉险滩、闯难关的深刻革命。

（2）它是一场思想理论的深刻变革——推动改革理论和政策的一系列新的重大突破。

（3）它是一场改革组织方式的深刻变革——加强党中央对重大工作的领导体制，完善党中央重大决策部署落实机制。

（4）它是国家制度和治理体系的深刻变革——在多年改革开放基础上的全面、系统、整体的制度创新，使社会主义制度的优越性得到更好发挥。

（5）它是人民广泛参与的深刻变革——尊重人民首创精神，激发人民群众的创造活力。

\noindent（二）以新发展理念引领高质量发展（5 分）：

（1）创新是引领发展的第一动力——坚持创新在我国现代化建设全局中的核心地位，加快实现高水平科技自立自强。

（2）协调是持续健康发展的内在要求——正确处理局部和全局、当前和长远、重点和非重点的关系，增强发展的整体性。

（3）绿色是永续发展的必要条件——实现经济社会发展和生态环境保护协同共进，加快发展方式绿色转型。

（4）开放是国家繁荣发展的必由之路——推动形成更大范围、更宽领域、更深层次对外开放格局，增强国际经济合作和竞争新优势。

（5）共享是中国特色社会主义的本质要求——坚持全民共享、全面共享、共建共享、渐进共享，不断推进全体人民共同富裕。

\vspace{4pt}
\noindent 综合来看，新发展理念是一个系统的理论体系，必须完整、准确、全面理解和把握，以引领高质量发展，为全面建设社会主义现代化国家提供更为坚实的物质基础。

\end{answerbox}

\vspace{12pt}
\examrule
\begin{center}
  \textit{参考答案仅供参考，评分时应酌情给分；要点正确、论述合理即可得满分或过程分。}
\end{center}

\end{document}